\appendices
\crefalias{section}{appendix}
\crefalias{subsection}{appendix}
\section{Operation Counts of the Precoders and Combiners}
\label{app:complexity}

This appendix derives the entries of \Cref{tab:complexity}. Operations are
counted as in \cite{peschiera2026fr3}: a complex multiply-accumulate is two
operations, so an inner product of two length-$r$ complex vectors costs $2r$
operations and the efficiencies $\eta$ of \eqref{eq:pmimo} are expressed per
operation. All counts below are per subcarrier and per coherence block; dividing
by $\upsilon_{\mathrm{coh}}$ turns them into the per-sample figures used in
\eqref{eq:gops_split}. The subcarrier index is dropped throughout.

\subsection{Computing the precoder or combiner}
Every scheme of \Cref{ssec:operation} other than \gls{mr} has the same three
stages, and the two families differ only in which dimension is inverted. Let
$\mat{A}\inC{n\times n}$ be the Hermitian positive definite matrix that is
factorized and let $r$ be the length of the vectors from which its entries are
built, which is also the number of right-hand sides that are solved against it.
The centralized precoder \eqref{eq:zf-precoder} and its uplink counterpart form
$\mat{A}=\mat{H}^{H}\mat{H}$ and solve $\bar{\mat{W}}\mat{A}=\mat{H}$ for the
$LM$ rows of $\mat{H}\inC{LM\times K}$, so $n=K$ and $r=LM$. The local schemes
\eqref{eq:local-rzf} and \eqref{eq:ul-local-mmse} form
$\mat{A}=\mat{E}_{l}\mat{E}_{l}^{H}+\lambda\mat{I}_{M}$ and solve
$\mat{A}\mat{W}_{l}=\mat{E}_{l}$ for the $K$ columns of
$\mat{E}_{l}\inC{M\times K}$, so $n=M$ and $r=K$. The three stages cost
\begin{enumerate}
\item \emph{Gram matrix.} $\mat{A}$ is Hermitian, so only its $n(n+1)/2$
upper-triangular entries are formed, each an inner product of length $r$ at $2r$
operations, giving $(n^{2}+n)\,r$. The diagonal loading $\lambda\mat{I}_{n}$ adds
$n$ real additions, and the power weighting $\mat{P}$ of
\eqref{eq:ul-local-mmse} rescales the $K$ columns of $\mat{E}_{l}$ beforehand at
$O(MK)$; both are dropped as lower order.
\item \emph{Cholesky factorization.} $\mat{A}=\mat{L}\mat{L}^{H}$ costs
$n^{3}/3$ operations, the count used in \cite{peschiera2026fr3}.
\item \emph{Triangular solves.} Forward and backward substitution against
$\mat{L}$ costs $n^{2}$ complex multiply-accumulates per right-hand side
\cite[App.~B.1.1]{massivemimobook}, hence $2n^{2}$ operations, and there are
$r$ of them, giving $2n^{2}r$.
\end{enumerate}
Summing the three stages,
\begin{equation}
\label{eq:app_chol}
\upsilon_{\mathrm{coh}}\,\Xi_{\mathrm{comp}}(n,r) =
\frac{n^{3}}{3} + 3n^{2}r + n\,r ,
\end{equation}
which is the second row of \Cref{tab:complexity} at $(n,r)=(K,LM)$ and the third
at $(n,r)=(M,K)$. Setting $r$ to the number of \gls{rf} chains of a single site
recovers the co-located count of \cite{peschiera2026fr3} exactly, so the
distributed and co-located deployments are priced by the same expression and can
be compared. The \gls{mr} row costs nothing to compute: under the perfect
\gls{csi} assumed here the direction is the channel itself,
$\bar{\vect{w}}_{k}=\vect{h}_{k}$ in the downlink and
$\vect{v}_{kl}=\vect{h}_{kl}$ in the uplink, so no factorization is performed.

Only the middle term of \eqref{eq:app_chol} is a genuine cubic cost in the
antenna count, and it is the one that separates the two families. At
$(n,r)=(K,LM)$ the factorization grows as $K^{2}LM$ and is carried by a single
node, whereas at $(n,r)=(M,K)$ each of the $L$ \glspl{ap} pays $M^{2}K$, for a
network total of $LM^{2}K$. Local processing is therefore cheaper than
centralized processing by a factor $K/M$ in this term, which is the
computational counterpart of the rate loss it accepts.

\begin{remark}[Convention on the cubic term]
\label{rem:cubic}
Counting a complex multiply-accumulate as two operations, as stages 1 and 3 do,
would give $2(n^{3}-n)/3$ rather than $n^{3}/3$ for the Cholesky factorization
\cite[App.~B.1.1]{massivemimobook}. We keep $n^{3}/3$ so that
\eqref{eq:app_chol} reproduces \cite{peschiera2026fr3} term for term and the
co-located baseline is not priced by a different rule than the distributed
network. The residual ambiguity is the ratio $n/(9r)$ of the factorization to
the dominant $3n^{2}r$ term, which is small wherever the matrix is small
relative to the number of right-hand sides. Both families satisfy this at the
operating point used here, $K=20$ \glspl{ue} against $LM=64$ antennas
centrally and $M=4$ antennas against $K=20$ \glspl{ue} locally, giving
$\SI{3.5}{\percent}$ and
$\SI{2.2}{\percent}$. It ceases to hold for the local schemes once an \gls{ap}
carries as many antennas as the network serves \glspl{ue}, so a deployment with
few large \glspl{ap} should be checked against the doubled term.
\todo{fix these numbers once $K$ and the deployment sweep are final.}
\end{remark}

\subsection{Applying the precoder or combiner}
Applying is scheme-independent. \gls{ap} $l$ forms
$\vect{x}_{l}=\mat{W}_{l}\vect{s}$ with $\mat{W}_{l}\inC{M\times K}$, which is
$MK$ complex multiplications and $M(K-1)$ complex additions, hence $2KM$
operations per sample to leading order. In the uplink the partial sums
$\vect{v}_{kl}^{H}\vect{y}_{l}$ of \eqref{eq:ul-distributive} cost the same,
$K$ inner products of length $M$. A node that processes all $LM$ antennas, which
is the \gls{cpu} under S1, pays $2K\,LM$ instead. This is the $2K$ operations
per sample and per processed antenna quoted in \Cref{sssec:mimo} and it appears
once per direction, since the precoded samples and the combined samples are
distinct data.

\subsection{What is not counted}
The power control of \eqref{eq:dl-fractional} and \eqref{eq:local-fractional}
needs the precoder norms $\sum_{q}\norm{\bar{\vect{w}}_{k}[q]}^{2}$ and the
large-scale coefficients \eqref{eq:beta}, which is $O(K\,LM)$ per coherence block
across the network and therefore below the $O(K^{2}LM)$ factorization it is
attached to. Channel estimation is not priced either, since perfect \gls{csi} is
assumed, and \Cref{tab:split} lists it for completeness only. Neither omission
is specific to the distributed deployment, so both cancel in the comparison of
\Cref{ssec:ee}.
